\chapter{Introducción}

\section{Antecedentes}

La búsqueda de vivienda en renta por parte de estudiantes universitarios que provienen de otras localidades distintas a Miahuatlán de Porfirio Díaz es un problema que ha sido abordado desde distintas perspectivas: el estudio del impacto de las universidades en el desarrollo local, el análisis del crecimiento urbano, el efecto de las plataformas digitales en el mercado de alquiler y el desarrollo de sistemas informáticos para la gestión de arrendamientos. En esta sección se presenta, en primer lugar, el panorama estadístico y territorial que enmarca el problema y, posteriormente, una síntesis crítica de las investigaciones previas más relevantes, con el propósito de identificar lo que ya se conoce, las limitaciones de las propuestas existentes y la brecha que el presente proyecto busca atender.

\subsection{Panorama estadístico y territorial}

De acuerdo con la Encuesta Nacional de Vivienda (ENVI) 2020, levantada mediante un muestreo probabilístico, bietápico, estratificado y por conglomerados de 55\,147 viviendas, de los 35.3 millones de viviendas particulares habitadas en México, el 16.4\,\% son rentadas; de estas 5.8 millones de viviendas en renta, únicamente el 54.0\,\% cuenta con un contrato de renta vigente \parencite{INEGI2021ENVI}. Es decir, casi la mitad de los arrendamientos del país se realiza sin un instrumento formal que dé certeza a arrendadores y arrendatarios. Asimismo, de acuerdo con el \textcite{INEGI2021ENVI}, el 21.1\,\% de los hogares (7.6 millones) reportó que alguno de sus integrantes necesita o planea rentar, comprar o construir una vivienda. Aunque la ENVI tiene representatividad nacional y estatal, no ofrece información a nivel municipal ni distingue la demanda específica de la población estudiantil.

En el ámbito local, la información censal del \textcite{INEGI2020DataMexico}, difundida a través de la plataforma Data México, describe el perfil habitacional de Miahuatlán de Porfirio Díaz: viviendas particulares habitadas, promedio de ocupantes por vivienda, distribución de habitaciones y acceso a servicios como internet, agua y electricidad.
% TODO: incorporar aquí las cifras concretas de Data México (viviendas habitadas, ocupantes por vivienda, % con internet).
Estos datos corresponden al corte censal de 2020 y no reflejan la oferta privada de renta ni la ubicación de los inmuebles disponibles. De forma complementaria, el Plan de Ordenamiento Territorial elaborado por el \textcite{SEINFRA2024} delimita las colonias del municipio, la cobertura de las redes de servicios básicos y la proximidad a nodos educativos como la Universidad de la Sierra Sur (UNSIS), con base en un diagnóstico territorial participativo y análisis cartográfico. Al ser un instrumento de planeación estático, no permite consultar la oferta de vivienda en renta, pero proporciona la base geográfica oficial para categorizar las zonas en las que se ubican los inmuebles.

\subsection{La universidad como generadora de demanda habitacional en Miahuatlán de Porfirio Díaz}

\textcite{ZepedaGarcia2019} examinaron cómo la llegada de la UNSIS ha transformado a Miahuatlán de Porfirio Díaz en un periodo aproximado de quince años, con el objetivo de valorar a la universidad como productora de capital humano y como impulsora de la economía local. Mediante una investigación de campo que recuperó testimonios de habitantes que presenciaron dichos cambios, complementada con análisis documental, las autoras encontraron que la llegada de estudiantes, profesores y personal administrativo modificó las condiciones de vida de la ciudad y generó una demanda creciente de servicios, entre ellos el alojamiento. Sus resultados, sin embargo, son de carácter cualitativo: no cuantifican la demanda de vivienda ni proponen mecanismos para organizarla. Este trabajo es relevante para el presente proyecto porque documenta, desde la investigación académica y en el mismo contexto geográfico, que la comunidad universitaria constituye un mercado permanente de demanda de vivienda en renta.

En una línea complementaria, \textcite{AvilaFlores2020} analizaron los factores del crecimiento urbano y del desarrollo sustentable en Miahuatlán de Porfirio Díaz durante el siglo \textsc{xxi}. A través de un análisis territorial, espacial y demográfico apoyado en Sistemas de Información Geográfica (SIG) y cartografía, estimaron la tasa de expansión urbana y la densidad de población, lo que les permitió identificar las zonas de mayor presión y expansión dentro del municipio. Su enfoque es macroscópico, por lo que no llega al nivel de la oferta inmobiliaria individual. Para esta investigación, dicho estudio aporta elementos para fundamentar el componente de geolocalización de la plataforma, al permitir identificar las zonas de expansión en las que resulta más necesario vincular la oferta de cuartos con los estudiantes.
% TODO: incorporar una cifra concreta de este estudio (tasa de expansión urbana o densidad) tras revisar el texto completo.

\subsection{Las plataformas digitales en el mercado de vivienda en alquiler}

De acuerdo a un análisis de distribución geográfica de los distintos tipos de alojamiento ofertados en Airbnb en la Ciudad de México entre 2019 y 2023, así como sus precios y su relación con el parque habitacional tradicional, con datos de la plataforma InsideAirbnb y del Censo de Población y Vivienda, la pandemia por COVID-19 redujo la oferta de alojamiento temporal en 2021 sin que los precios disminuyeran. Posteriormente la plataforma se recuperó hasta superar la expansión observada a inicios de 2019, en algunos territorios más del 60\,\% del parque habitacional se encuentra inserto en la plataforma, con concentraciones en las alcaldías Miguel Hidalgo y Cuajimalpa. La propia autora reconoce que la evidencia es insuficiente para determinar los efectos de esta forma de alquiler en cada municipio y que en México la investigación sobre el fenómeno es escasa; además, su análisis se centra en el alquiler turístico de corta estancia en una metrópoli \parencite{GarridoRodriguez2025}. Para el presente proyecto, este trabajo muestra que las plataformas digitales reorganizan efectivamente el mercado de alquiler y que el análisis geográfico de la oferta es una herramienta útil para comprenderlo, aunque no existen estudios equivalentes para el alquiler de mediana estancia dirigido a estudiantes en municipios pequeños.

\subsection{Sistemas web para la gestión de arrendamientos}

El antecedente tecnológico más cercano a esta investigación es el de \textcite{MejiaAvecias2023}, quienes desarrollaron un sistema web para facilitar la obtención y distribución de información sobre departamentos en renta para estudiantes foráneos de la Escuela Superior de Tlahuelilpan, Hidalgo. A partir de información recabada con estudiantes de la zona, identificaron que la oferta se difundía principalmente por redes sociales, con información escasa, ambigua y dispersa, e incluso con publicaciones de más de un año de antigüedad sin actualizar. El desarrollo siguió la metodología ágil Scrum; los requerimientos se obtuvieron mediante entrevistas a estudiantes foráneos y el sistema se construyó con una arquitectura cliente-servidor empleando HTML, CSS, JavaScript, PHP y MySQL. Como resultado, presentaron el modelo de datos, el diagrama de módulos y doce interfaces que incluyen el registro de inmuebles, la búsqueda con filtros por ubicación, tamaño, servicios y costo, la reservación en línea y la comunicación entre arrendadores y estudiantes. Asimismo, al revisar plataformas existentes como Stukey, Roomgo o Spotahome, señalan que estas se concentran en la Ciudad de México y en otras ciudades grandes del país. No obstante, el trabajo no reporta una evaluación empírica del sistema, como pruebas con usuarios, métricas de usabilidad o tiempos de búsqueda, y fue diseñado para el contexto de Tlahuelilpan. Este antecedente sirve como guía para definir el modelo de datos, los módulos y los filtros de búsqueda de la plataforma propuesta, la cual, a diferencia de este, contemplará la evaluación con usuarios y las condiciones propias de Miahuatlán de Porfirio Díaz.

Por su parte, \textcite{Valenzuela2024} desarrolló un sistema web para optimizar el control de pagos por arrendamiento de locales en la Sub Gerencia de Recaudación de la Municipalidad Provincial de Pasco, Perú, ante la deficiencia, lentitud y falta de trazabilidad del control manual de dichos pagos. Mediante un método deductivo e inductivo, con una población de 100 arrendatarios y una muestra de 20, evaluó los tiempos de procesamiento de pagos, la reducción de errores de registro y la generación de reportes, y concluyó que la implementación mejoró el control, la rapidez y la confiabilidad de los arrendamientos.
% TODO: incorporar los porcentajes de mejora que reporta la tesis tras revisar el documento completo.
Sin embargo, su lógica de negocio está orientada a locales comerciales administrados por una institución municipal y no al alojamiento estudiantil. Para el presente proyecto, este trabajo ofrece referencias técnicas sobre el modelado de la base de datos y de los flujos necesarios para gestionar contratos y pagos dentro de un sistema de rentas.

\subsection{Síntesis y brecha identificada}
Los estudios sobre Miahuatlán de Porfirio Díaz \parencite{AvilaFlores2020, ZepedaGarcia2019} demuestran que la presencia de la Universidad de la Sierra Sur (UNSIS) genera una demanda habitacional constante y que existen herramientas de análisis territorial para ubicarla; sin embargo, ninguno propone una solución que organice la oferta de vivienda. La literatura sobre plataformas digitales \parencite{GarridoRodriguez2025} muestra que estas reorganizan el mercado de alquiler, pero se ha concentrado en el alquiler turístico de las grandes ciudades. Finalmente, los sistemas web existentes \parencite{MejiaAvecias2023, Valenzuela2024} prueban la viabilidad técnica de centralizar la información de arrendamiento, pero fueron diseñados para otros contextos o tipos de inmueble y no reportan una evaluación con usuarios finales en el caso de la vivienda estudiantil.

En consecuencia, no se identificó una solución que integre la oferta de cuartos y departamentos en renta de Miahuatlán de Porfirio Díaz con información geográfica sobre colonias, servicios y cercanía a la UNSIS, dirigida a personas no residentes y evaluada con sus usuarios. Esta es la brecha que el presente proyecto busca atender.

\section{Planteamiento del problema}

Miahuatlán de Porfirio Díaz, Oaxaca, se ha consolidado progresivamente como un importante centro comercial y académico de la región, lo que ha propiciado la llegada de personas provenientes de distintas localidades que se trasladan a este municipio por motivos laborales y educativos. Entre este sector de la población destacan los estudiantes de la Universidad de la Sierra Sur (UNSIS), cuya matrícula es de aproximadamente 2,294 estudiantes, de los cuales un porcentaje considerable requiere trasladarse temporalmente a Miahuatlán para continuar con su formación académica. Esta situación genera una demanda constante de viviendas en renta, principalmente entre quienes no cuentan con un lugar de residencia permanente en la localidad.

Actualmente, la difusión de la oferta de viviendas en renta se realiza principalmente mediante grupos de redes sociales, como Facebook, así como a través de anuncios físicos colocados en postes y espacios cercanos a las instalaciones de la universidad. Sin embargo, estos medios funcionan de manera independiente y no existe un canal centralizado que permita consultar, organizar y actualizar la información de los inmuebles disponibles. Esta dispersión dificulta el acceso a información confiable y oportuna, ya que las publicaciones pueden quedar desactualizadas, presentar datos incompletos o carecer de elementos que permitan verificar la identidad de los arrendadores y las condiciones reales de las viviendas ofertadas.

Como consecuencia, las personas que buscan alojamiento, particularmente los estudiantes no residentes, enfrentan dificultades para localizar y comparar alternativas de vivienda de acuerdo con su presupuesto, ubicación, características y necesidades particulares. Esto incrementa el tiempo y esfuerzo dedicado al proceso de búsqueda, además de generar incertidumbre al momento de establecer contacto con los arrendadores. Asimismo, la falta de mecanismos de verificación y de información confiable aumenta el riesgo de ser víctimas de posibles fraudes, como ofertas de inmuebles inexistentes, publicaciones engañosas o viviendas cuyas condiciones no corresponden con lo anunciado, lo que puede ocasionar pérdidas económicas y afectar la experiencia de adaptación de los estudiantes a su nuevo entorno.

En este contexto, se identifica la necesidad de contar con una solución tecnológica que centralice la oferta de viviendas en renta en Miahuatlán de Porfirio Díaz, facilite la búsqueda y comparación de inmuebles mediante criterios específicos y proporcione información organizada y actualizada. Asimismo, resulta importante incorporar mecanismos que contribuyan a generar mayor confianza entre arrendadores y arrendatarios, favoreciendo un proceso de búsqueda y selección de vivienda más accesible, eficiente y seguro, especialmente para la comunidad estudiantil de la Universidad de la Sierra Sur.
